\begin{table}[!htbp]\centering\scriptsize
\setlength{\tabcolsep}{3.5pt}
\caption{Transition année par année, avec les bornes (Md€ courants)}\label{tab:is_transition_complet}
\begin{tabular}{@{}l>{\raggedleft\arraybackslash}p{1.3cm}>{\raggedleft\arraybackslash}p{1.3cm}>{\raggedleft\arraybackslash}p{1.3cm}>{\raggedleft\arraybackslash}p{1.3cm}>{\raggedleft\arraybackslash}p{1.3cm}>{\raggedleft\arraybackslash}p{1.3cm}>{\raggedleft\arraybackslash}p{1.3cm}>{\raggedleft\arraybackslash}p{1.3cm}@{}}
\toprule
& \multicolumn{4}{c}{Cas central} & \multicolumn{2}{c}{Cas favorable} & \multicolumn{2}{c}{Cas défavorable}\\
\cmidrule(lr){2-5}\cmidrule(lr){6-7}\cmidrule(lr){8-9}
Année & Coût & Ressources & Solde & Dette & Coût & Dette & Coût & Dette \\
\midrule
2027 & 8,9 & 4,1 & $-$4,8 & 4,8 & 8,4 & 3,8 & 9,3 & 5,8 \\
2028 & 8,3 & 4,2 & $-$4,1 & 9,1 & 7,1 & 6,3 & 9,4 & 11,8 \\
2029 & 7,9 & 4,3 & $-$3,6 & 12,9 & 6,3 & 7,9 & 9,4 & 17,9 \\
2030 & 7,5 & 4,4 & $-$3,1 & 16,4 & 5,9 & 8,9 & 9,2 & 23,9 \\
2031 & 7,2 & 4,6 & $-$2,6 & 19,6 & 5,6 & 9,6 & 8,8 & 29,6 \\
2032 & 6,8 & 4,7 & $-$2,0 & 22,3 & 5,2 & 9,8 & 8,3 & 34,8 \\
2033 & 4,7 & 4,8 & 0,1 & 22,9 & 3,4 & 7,9 & 6,1 & 37,8 \\
2034 & 3,6 & 5,0 & 1,4 & 22,2 & 2,6 & 5,1 & 4,5 & 39,2 \\
2035 & 2,9 & 5,1 & 2,2 & 20,7 & 2,1 & 1,6 & 3,7 & 39,8 \\
2036 & 2,4 & 5,3 & 2,9 & 18,5 & 1,8 & $-$2,6 & 3,0 & 39,5 \\
2037 & 2,1 & 5,5 & 3,4 & 15,7 & 1,6 & $-$7,3 & 2,6 & 38,7 \\
2038 & 1,9 & 5,6 & 3,7 & 12,5 & 1,5 & $-$12,5 & 2,3 & 37,5 \\
2039 & 1,9 & 5,8 & 3,9 & 9,1 & 1,5 & $-$17,9 & 2,3 & 36,0 \\
2040 & 2,0 & 6,0 & 4,0 & 5,4 & 1,6 & $-$23,7 & 2,4 & 34,4 \\
2041 & 2,0 & 6,1 & 4,1 & 1,4 & 1,6 & $-$29,8 & 2,4 & 32,7 \\
2042 & 2,1 & 6,3 & 4,2 & $-$2,7 & 1,7 & $-$36,3 & 2,5 & 30,8 \\
2043 & 2,2 & 6,5 & 4,3 & $-$7,2 & 1,7 & $-$43,1 & 2,6 & 28,8 \\
2044 & 2,2 & 6,7 & 4,5 & $-$11,9 & 1,8 & $-$50,4 & 2,7 & 26,6 \\
2045 & 2,3 & 6,9 & 4,6 & $-$16,9 & 1,8 & $-$58,1 & 2,7 & 24,3 \\
2046 & 2,4 & 7,1 & 4,8 & $-$22,2 & 1,9 & $-$66,2 & 2,8 & 21,8 \\
2047 & 2,4 & 7,3 & 4,9 & $-$27,8 & 2,0 & $-$74,7 & 2,9 & 19,1 \\
2048 & 2,5 & 7,6 & 5,0 & $-$33,8 & 2,0 & $-$83,7 & 3,0 & 16,2 \\
2049 & 2,6 & 7,8 & 5,2 & $-$40,0 & 2,1 & $-$93,2 & 3,1 & 13,1 \\
2050 & 2,7 & 8,0 & 5,3 & $-$46,7 & 2,1 & $-$103,2 & 3,2 & 9,8 \\
2051 & 2,7 & 8,3 & 5,5 & $-$53,7 & 2,2 & $-$113,8 & 3,3 & 6,3 \\
2052 & 2,8 & 8,5 & 5,7 & $-$61,1 & 2,3 & $-$124,9 & 3,4 & 2,6 \\
2053 & 2,9 & 8,8 & 5,8 & $-$69,0 & 2,3 & $-$136,6 & 3,5 & $-$1,4 \\
2054 & 3,0 & 9,0 & 6,0 & $-$77,2 & 2,4 & $-$148,9 & 3,6 & $-$5,6 \\
2055 & 3,1 & 9,3 & 6,2 & $-$86,0 & 2,5 & $-$161,8 & 3,7 & $-$10,1 \\
2056 & 3,2 & 9,6 & 6,4 & $-$95,1 & 2,6 & $-$175,4 & 3,8 & $-$14,8 \\
\bottomrule
\end{tabular}
\par\vspace{3pt}\parbox{\linewidth}{\raggedright\scriptsize Cas favorable : régime actuel dégressif, ressources hautes (4,62~Md€ en 2027). Cas défavorable : régime actuel linéaire, ressources basses (3,50~Md€). Cas central : moitié-moitié, ressources au milieu (4,06~Md€). Une dette négative est un excédent cumulé : dans la proposition, il n'apparaît pas, car l'excédent est rendu par la baisse du taux normal dès l'année qui suit l'extinction (le tableau montre ce qui se passerait sans cette baisse). Part imposable : 75~\%. Source : \texttt{is/calculs\_note\_is.py}.}
\end{table}
